\SourceLeaf{443}%
Zusammentreffen zweier Linien gebildet werden: offenbar bleiben daher auch solche Linien, die auf der einen Fläche kürzeste Linien waren, bei der Übertragung wieder kürzeste Linien. Hieraus folgt leicht, dass wenn einem beliebigen durch kürzeste Linien gebildeten Polygon, in so fern es auf der ersten Fläche ist, die Figur der Zenithe auf der Hülfskugel entspricht, deren Area $=A$ ist, demselben Polygon hingegen, in so fern es durch Abwickelung auf die zweite Fläche übertragen wird, eine Figur der Zenithe auf der Hülfskugel, deren Inhalt $=A'$ wird, allemal
\[
A=A'
\]
sein wird. Obgleich dieser Beweis ursprünglich die Begrenzung der Figuren durch kürzeste Linien voraussetzt, so sieht man doch leicht, dass er allgemein gültig ist, wie auch die Begrenzung sein möge: denn in der That, wenn das Theorem von der Anzahl der Seiten unabhängig ist, so hindert nichts, für jedes Polygon, dessen Seiten alle oder einige keine kürzeste Linien sind, ein anderes von unendlich vielen Seiten, die jede kürzeste Linien sind, zu denken.

Ferner ist klar, dass bei der Übertragung durch Abwickelung jede Figur auch ihre Area behält.

Es ist nun hier von $4$ Figuren die Rede:
\begin{enumerate}
\item[1)] eine beliebige Figur auf der ersten Fläche,
\item[2)] die Figur auf der Hülfskugel, die den Zenithen von jener entspricht,
\item[3)] die Figur auf der zweiten Fläche, welche Nro $1$ vermöge der Abwickelung bildet,
\item[4)] die Figur auf der Hülfskugel, die den Zenithen von Nro $3$ entspricht.
\end{enumerate}
\noindent\hspace*{1.2em}Nach dem, was wir bewiesen haben, haben also $2$ und $4$ gleichen Inhalt, eben so wie $1$ und $3$. In so fern nun diese Figuren unendlich klein angenommen werden, ist der Quotient, wenn man $2$ durch $1$ dividirt, das Krümmungsmass an der Stelle der ersten krummen Fläche, und eben so der Quotient $4$ durch $3$ bei der zweiten. Hieraus folgt also der wichtige Lehrsatz, dass
\begin{quote}
»bei der Übertragung der Flächen durch Abwickelung das Krümmungsmass an jeder Stelle unverändert bleibt«.
\end{quote}
\noindent Dasselbe gilt also von dem Product des grössten und kleinsten Krümmungshalbmessers.

Bei der Ebene ist offenbar das Krümmungsmass überall $=0$: hieraus
